\chapter{Slow Heat}


% -----------------------------------------------------------
\section{SlowHEAT -- Heat the Body, Not the Building}
\textit{Synthesis of the presentation by Prof.\ Geoffrey van Moeseke,
UCLouvain, February 18, 2026}

% -----------------------------------------------------------
\subsection{Core Intuition and Motivation}

The buildings sector is one of the largest contributors to
greenhouse gas emissions in Europe, and space heating accounts for
the bulk of residential energy consumption. The dominant policy
response has been to invest massively in building renovation ---
improving insulation, replacing windows, and upgrading heating
systems. The EU renovation wave is estimated to require roughly
\texteuro243 billion per year to achieve its climate objectives.
Yet a simple thought experiment reveals the staggering scale of
this commitment: equipping every EU citizen with high-quality
insulating outerwear would cost approximately \texteuro12.5 billion
per year --- nearly twenty times less.

This provocative comparison is not a call to replace renovation
with ski suits, but rather an invitation to question a deeply
embedded assumption: that thermal comfort necessarily means heating
the entire volume of a building to a uniform temperature. The
SlowHEAT project, led by Prof.\ Geoffrey van Moeseke at UCLouvain
and conducted in collaboration with UCLouvain's iPSY institute,
the \textit{habitat et participation} association, LAB Architecture,
\textit{Architecture et Climat}, and ULB-IGEAT, with support from
Innoviris Brussels, explores a radically different paradigm.
Rather than heating the building, the core idea is to direct heat
toward the occupant's body through \textit{Personal Comfort Systems}
(PCS) --- a concept with deep historical roots and a growing body
of scientific support.

% -----------------------------------------------------------
\subsection{An Idea Rooted in History}

The notion of localised, body-centred heating is far from new.
For millennia, heat was confined to specific spots within the
dwelling. Ancient Greek braziers, Gaulish central hearths, Viking
longhouse fire pits --- all provided radiant warmth to those
gathered around them, while leaving the rest of the space cold.
Furniture was similarly designed to concentrate warmth: canopied
and curtained beds created insulated sleeping micro-environments,
hooded chairs shielded occupants from draughts, folding screens
blocked cold air, and small portable foot warmers carried embers
beneath one's skirts. Japanese kotatsu tables, combining a low
table with a heating element and a heavy draping blanket, are a
particularly elegant example of the same principle: heat the body,
not the room.

The shift toward uniform centralised heating began in earnest in
the late nineteenth and early twentieth centuries, with the
industrialisation of cast-iron radiators and gas or coal boilers.
This transition was not purely technological --- it was also
cultural. Advertising campaigns of the 1930s (including those of
the Compagnie Nationale des Radiateurs in France and the National
Radiator Company in England) actively constructed a new narrative:
central heating was modern, hygienic, uniform, effortless, and
liberating. The old way of warming --- gathering around a stove,
layering clothing indoors, retreating to a heated corner of the
house --- was reframed as backward, uncomfortable, and redolent
of poverty.

The consequences have been significant. Historical data show that
average indoor winter temperatures in UK homes rose from
approximately 15--16\,\textdegree C in the 1960s to
19--20\,\textdegree C by the late 1990s. Research by Mavrogianni
et al.\ (2013) documents this steady upward drift across living
rooms, bedrooms, and hallways alike. At the EU scale, the European
Environment Agency has estimated that switching from single-room
heating to central heating increases energy demand for space
heating by roughly 25\,\% on average --- an increase that has
partly offset the efficiency gains achieved through better
insulation and more efficient boilers over the same period.

% -----------------------------------------------------------
\subsection{An Idea Aiming for Progress}

Far from being a nostalgic regression, the SlowHEAT approach is
positioned as a forward-looking response to multiple contemporary
challenges.

\paragraph{Improved health.}
Recent work published in \textit{Energy \& Buildings} (Zhu et al.,
2025) underscores that prolonged exposure to thermoneutral
environments may compromise both comfort and health. The human
body is equipped with a sophisticated array of thermoregulatory
mechanisms --- vasoconstriction and vasodilation, piloerection,
brown adipose tissue activation, and shivering thermogenesis ---
that are simply not engaged when indoor temperatures are
maintained at a constant 20--22\,\textdegree C. Mild, repeated
cold exposure has been associated with metabolic benefits,
including activation of brown fat, as demonstrated by van der
Lans et al.\ (2013). The hormesis principle --- whereby moderate
physiological stress produces adaptive benefits --- may apply to
thermal environments just as it does to exercise.

\paragraph{Improved individual wellbeing.}
Unlike conventional HVAC systems that impose a single thermal
setpoint on all occupants of a shared space, PCS allow individuals
to modulate their own thermal environment without affecting those
around them. Research on office environments (Wegertseder-Martinez
et al.) confirms that this individualised control improves both
comfort satisfaction and perceived wellbeing.

\paragraph{A pathway toward electrification.}
The phase-out of fossil fuel heating is a central objective of
European climate policy. In Belgium, oil boilers are already
banned in new constructions in Wallonia (from January 2026), and
Brussels has extended the ban to gas boilers as well. Heat pumps
are the primary proposed alternative, but their high upfront cost
and the need for low-temperature distribution systems make them
challenging to retrofit. Low-power electric PCS --- consuming
tens to a few hundred watts rather than several kilowatts ---
represent a more immediately accessible bridge technology,
compatible with existing electrical infrastructure.

\paragraph{Technological and social pluralism.}
The approach does not prescribe a single solution. It
accommodates both high-tech options (battery-powered heated
garments, smart infrared panels, app-controlled electric blankets)
and low-tech ones (well-chosen woollen underlayers, homemade
heated capes fabricated in community repair workshops, draught
exclusion). This breadth means it can appeal to both
techno-optimists and advocates of degrowth and sufficiency.

% -----------------------------------------------------------
\subsection{The SlowHEAT Living-Lab Experiment (2020--2023)}

\subsubsection{General Framework}

The empirical core of SlowHEAT is a three-year living-lab study
conducted with thirty volunteer households in Belgium, running from
winter 2020--2021 through winter 2022--2023. The study period
coincided, unplanned, with the 2021--2022 energy price crisis
triggered by the Russian invasion of Ukraine --- an external shock
that adds both complexity and ecological validity to the results.

Participants were voluntary, self-selected, and environmentally
motivated --- a sample the authors themselves acknowledge as likely
to produce over-optimistic results compared to the general
population. The group was predominantly highly educated (64\,\%
held a master's degree, 14\,\% a PhD), with a balanced gender
distribution and a spread of age groups from the 20s to the 60s,
with the largest cohort in their 40s (31\,\%). Housing types and
energy performance levels were deliberately varied.

The methodological framework drew on social practice theory,
following Schatzki (1997) and Gram-Hanssen (2010), analysing
heating behaviour along four interacting dimensions: material
infrastructures and objects, practical knowledge and skills,
representations and meanings, and collective norms and rules.
Participants were not asked to endure discomfort: the goal was to
find a personal equilibrium, supported by a toolkit of PCS made
available to all households. These ranged from heated bed covers
(100\,W) and heating chair pads (30\,W) to small infrared panels
(300\,W), radiant heaters (1\,200\,W), large infrared panels
(400\,W), and heated desk pads (80\,W).

\subsubsection{Quantitative Results}

The central quantitative findings are striking. Over the course
of the project, participants progressively reduced their indoor
temperatures well below conventional norms. Table~\ref{tab:temps}
summarises the evolution of living-room temperatures.

\begin{table}[ht]
\centering
\caption{Indoor living-room temperature: thermostat setpoints
before the project and measured temperatures during the project.
Medians, means, and standard deviations are calculated for the
whole participant group. Both occupied and unoccupied periods
are included.}
\label{tab:temps}
\begin{tabular}{llccc}
  \toprule
  \textbf{Period} & \textbf{Metric} & \textbf{Median (\textdegree C)}
    & \textbf{Mean (\textdegree C)} & \textbf{SD (\textdegree C)} \\
  \midrule
  Before project & Thermostat setting & 19.0 & 19.0 & 1.6 \\
  \midrule
  \multirow{2}{*}{First winter}
    & Coldest week & 16.7 & 17.6 & 1.9 \\
    & Coldest day  & 15.8 & 16.9 & 2.1 \\
  \midrule
  \multirow{4}{*}{Third winter}
    & Coldest two months & 15.0 & 15.1 & 1.9 \\
    & Coldest month      & 14.5 & 14.7 & 2.0 \\
    & Coldest week       & 13.7 & 13.5 & 2.6 \\
    & Coldest day        & 12.5 & 12.6 & 2.5 \\
  \bottomrule
\end{tabular}
\end{table}

In terms of energy consumption, participants achieved average
reductions of approximately 50\,\% in gas consumption on a
heating-degree-day adjusted basis, and approximately 14\,\% in
electricity consumption. The authors note two important caveats:
gas savings may be overestimated because most dwellings are
apartments or semi-detached houses that benefit from heat
transfers from neighbouring units; and the electricity savings
were unexpected, likely reflecting the broader behavioural
response to the 2022 energy price shock rather than the PCS
alone.

Inter-household variation was considerable. Mean indoor
temperatures across participating dwellings ranged from
approximately 11\,\textdegree C to 16\,\textdegree C, reflecting
wide differences in housing performance, household composition,
and personal thermal preferences. One particularly illustrative
case is that of a household that had already undergone a building
renovation: over the same period, SlowHEAT practices generated
savings of 700\,m$^3$ of gas per year, compared to 380\,m$^3$
per year attributable to the renovation itself.

% -----------------------------------------------------------
\subsection{Qualitative Analysis: A Social-Practice Framework}

\subsubsection{Changing Representations}

Perhaps the deepest shift required by the SlowHEAT approach is
representational. Our current expectation that a home should be
uniformly heated to 19--22\,\textdegree C is not a biological
necessity but a historically and socially constructed norm ---
one that, as historian Olivier Jandot argues, is inseparable from
the technological evolution of the nineteenth and twentieth
centuries. This norm is recent, geographically variable, and
therefore changeable.

One of the most striking indicators of this normative rigidity
is the public reaction to media coverage of the SlowHEAT project.
Social media comments cited in the presentation include
accusations of endangering children, comparisons to medieval
living conditions, and assertions that living below 20\,\textdegree
C is simply unacceptable in 2022. These reactions illustrate the
depth of the cultural embedding of centralised warmth as a marker
of dignity and modernity --- an embedding that was itself
deliberately constructed by the heating industry a century ago.

Alongside this cultural shift, participants were invited to
develop a richer vocabulary for bodily thermal sensations and to
develop awareness of their body's own thermoregulatory capacities:
vasoconstriction, brown adipose tissue activation, and the
adaptive reduction in perceived cold discomfort with regular
mild exposure.

\subsubsection{Acquiring Knowledge and Skills}

A significant part of the learning process concerned the
relationship between indoor temperature, relative humidity, and
air quality. A key practical insight is that, at 15\,\textdegree
C, maintaining indoor relative humidity below approximately
55\,\% (corresponding to a humidity ratio of roughly
0.006\,g$_\text{water}$/kg$_\text{air}$) is sufficient to
prevent mould growth on most building materials. Critically, in
a Belgian winter climate, outdoor air humidity rarely exceeds
this threshold, meaning that regular ventilation with outdoor
air is almost always safe from a mould-prevention perspective ---
a counterintuitive finding for many participants who feared that
cooler interiors would be moister and more prone to mould.

The SlowHEAT team also developed a dedicated tool, the
\textit{SlowZone} widget, which combines real-time outdoor
temperature and humidity data with psychrometric calculations
to indicate the safe ventilation window for any given location.

Participants also developed practical skills in selecting and
deploying PCS: understanding the difference between conductive
heating (body contact), radiant heating (infrared proximity),
and convective heating (air warming), and learning which
solution best matched each activity, body zone, and time of day.

\subsubsection{Infrastructure and Technology: A Hierarchy of Solutions}

The project proposes a structured hierarchy of heating
interventions, ordered from least to most energy-intensive, as
summarised in Table~\ref{tab:hierarchy}.

\begin{table}[ht]
\centering
\caption{Proposed hierarchy of Personal Comfort System
interventions, from zero-energy measures to whole-building
heating.}
\label{tab:hierarchy}
\begin{tabular}{>{\bfseries}p{2.8cm} p{2.2cm} p{3.5cm} p{4cm}}
  \toprule
  Class & Power & Solution family & Examples \\
  \midrule
  A -- Zero-energy measures
    & 0\,W
    & Clothing, zoning, activity
    & Thermal underwear, closing doors, physical activity, curtains \\
  \addlinespace
  B -- Body-contact conduction
    & $\pm$50\,W/person
    & Wearable or contact accessories
    & Heated chair pad, electric blanket, heated cape, hot-water bottle \\
  \addlinespace
  C -- Radiant proximity
    & $\pm$300\,W/person
    & Radiant or heated furniture near occupant
    & Infrared panel, radiant column, heated table \\
  \addlinespace
  D -- Local room heating
    & $\pm$1\,500\,W/room
    & Thermostatic room heater
    & Room maintained at 15--17\,\textdegree C when occupied,
      with Classes A--C still active \\
  \addlinespace
  E -- Whole-building heating
    & $\pm$5\,000\,W/dwelling
    & Central heating system
    & Kept at frost-protection level (8\,\textdegree C) or
      passage temperature (12--15\,\textdegree C) \\
  \bottomrule
\end{tabular}
\end{table}

The philosophy is not to eliminate central heating entirely, but
to relegate it to a supporting, background role while Classes A
through C meet most comfort needs. Low-tech Class A and B
solutions are emphasised: the project ran community sewing
workshops to produce washable heated capes equipped with recycled
heating textile, and explored prototype heated chair inserts and
desk setups using low-voltage heating elements.

\subsubsection{Navigating Social Norms}

The social dimension proved one of the most challenging aspects
of the experiment. Participants navigated real tensions between
their new practices and the expectations of visitors, family
members, and social convention. Receiving guests in a
15\,\textdegree C home, asking in-laws to put on a jumper, or
negotiating thermostat settings within a household all require
a form of social courage and communication that purely
technological solutions do not demand.

Some participants experienced these tensions positively, as a
sign of being ahead of their time. Others found the social cost
significant. The project also surfaced a regulatory tension:
current energy performance certificate frameworks in Wallonia
(and more broadly across Europe) rate buildings on the basis
of normalised energy consumption at standard occupancy and
temperature assumptions. A household that genuinely heats its
home to 14\,\textdegree C and achieves excellent comfort through
PCS may appear, on paper, to have a poorly performing building
--- or may even find that certain renovation subsidies are
inaccessible because its actual consumption already falls below
the required thresholds. The authors argue that regulation should
distinguish more clearly between the \emph{means} (energy in
kWh/m$^2$) and the \emph{end} (a healthy, enabling domestic
environment), and that ``off-framework'' renovation pathways
tailored to sober lifestyles should be legally recognised.

% -----------------------------------------------------------
\subsection{Discussion: Efficiency, Sufficiency, and Beyond}

SlowHEAT raises fundamental questions about the relationship
between efficiency and sufficiency in energy policy. The dominant
paradigm --- enshrined in frameworks such as the Negawatt approach
--- places sufficiency first, followed by efficiency, then
renewable energy. In practice, however, sufficiency is rarely
operationalised with the same rigour as efficiency: it tends to
mean marginal reductions (lower the thermostat by 1\,\textdegree
C, reduce shower time by two minutes) rather than a qualitative
transformation of how services are delivered.

SlowHEAT illustrates what a more radical form of sufficiency can
look like in practice: not ``less of the same'' (a slightly cooler
uniformly heated home) but a fundamentally different service
configuration (a cool home with locally warm people). This
distinction is important not only for its energy implications but
for its social ones. As the concept of \textit{consumption
corridors} (Akenji et al., 2021) suggests, sufficiency should
be bounded both by a social floor --- below which underconsumption
causes deprivation and health risk --- and an environmental ceiling
--- above which overconsumption is ecologically unsustainable.
Both boundaries should be democratically established and informed
by science. The SlowHEAT experiment suggests that the social
floor for thermal comfort may be considerably lower than current
norms imply, provided adequate PCS are available and occupants
have the knowledge, agency, and social support to use them.

The project's main acknowledged limitations are twofold. First,
the sample's high educational level and prior environmental
motivation make it difficult to generalise: a participant group
with mixed motivations and lower climate awareness might achieve
smaller reductions and report lower satisfaction. Second, the
energy savings attributed to SlowHEAT may be partially
overestimated in multi-unit dwellings, where heat transfers from
neighbouring heated apartments contribute to maintaining indoor
temperatures. More controlled studies, with matched comparison
groups and sub-metered heating data, are needed to isolate the
specific contribution of PCS from confounding behavioural and
contextual factors.

Notwithstanding these limitations, the SlowHEAT project makes a
compelling empirical case that a significant fraction of European
households could, with appropriate support, substantially reduce
their heating energy consumption while maintaining subjective
comfort --- and that the barriers to doing so are as much
cultural, regulatory, and social as they are technical.